\documentclass[10pt,a4paper]{amsart}
\usepackage[margin=19mm]{geometry}
\usepackage{amsmath,amssymb,mathtools}
\usepackage[scr=rsfs,cal=euler]{mathalfa}
\usepackage{xfrac}
\usepackage[unicode,hidelinks,bookmarks=false]{hyperref}
\newcommand{\ie}{i.e.\ }
\newcommand{\cf}{cf.\ }
\newcommand{\resp}{respectively\ }
\newcommand{\Subsection}{\subsection}
\providecommand{\nobreakdash}{\nobreak\hbox{-}}
\setlength{\emergencystretch}{2em}
\multlinegap0pt
\usepackage{bm}
\usepackage{bbm}
\usepackage{dsfont}
\usepackage{xspace}
\usepackage{cancel}
\usepackage{mathtools}
\usepackage{amsthm}
\makeatletter
\@ifundefined{theorem}{\newtheorem{theorem}{Theorem}}{}
\@ifundefined{lemma}{\newtheorem{lemma}[theorem]{Lemma}}{}
\makeatother
\title[Posterior Concentration for Gaussian Process Priors under Re]{Posterior Concentration for Gaussian Process Priors under Rescaled and Hierarchical Mat\'ern and Confluent Hypergeometric Covariance Functions}
\author{Xiao Fang, Anindya Bhadra}
\date{Source: arXiv:2312.07502v3 (CC BY 4.0)}
\begin{document}
\maketitle

\noindent\emph{Source and licence.} Posterior Concentration for Gaussian Process Priors under Rescaled and Hierarchical Mat\'ern and Confluent Hypergeometric Covariance Functions. Version arXiv:2312.07502v3 (CC BY 4.0), distributed under the Creative Commons Attribution 4.0 International licence, as recorded in the arXiv licence metadata for this exact version and shown on the abstract page https://arxiv.org/abs/2312.07502v3. Full rights notice: licenses/arxiv-2312.07502-CC-BY-4.0.txt.\par\smallskip

\noindent\textit{Editorial scope.} Counted unit: the Lemma whose source label is \texttt{CH spectral density tail behavior} in the article (source0.tex lines 1137--1146, source label \texttt{CH spectral density tail behavior}), delivered together with the article's own complete original proof of it (source0.tex lines 1147--1198, one \texttt{proof} environment of 52 lines). No mathematics is written, repaired or re-derived: statement and proof are the article's own text line for line. Required context retained uncounted: incomplete gamma inequality (lines 1037--1039); incomplete gamma go to 1 (lines 1050--1052); incomplete integrate (lines 1083--1089). Every definition, display and label of those lines is retained unchanged. Omitted from this package and not redistributed: 1--1036, 1040--1049, 1053--1082, 1090--1136, 1199--1270. Nothing inside a retained excerpt is omitted or altered: every retained body is delivered exactly as the article's lines are written, so no omission receipt of any kind accompanies this package. Citations: the retained excerpts carry no citation command at all, so no bibliography entry of the article is transcribed and no reference is redistributed. Reference adaptation: none was needed and none was made; every reference inside the delivered unit resolves inside it, so no unresolved reference or citation is produced. Printed slips kept literal: no printed word, symbol, hypothesis or number of the counted statement or of its proof was altered. Layout and class: the article's own class is \texttt{imsart} with packages including amsthm, amsmath, amsfonts, amssymb, natbib, xcolor, hyperref, graphicx, bm, subfigure; the wrapper is the lane's pinned \texttt{amsart} editorial wrapper at 10pt on A4, which restates the theorem-like environments the retained text needs and adds the article's own macro definitions that the retained text needs (none), with extra wrapper packages (bm, bbm, dsfont, xspace, cancel, mathtools). The delivered numbering is the unit's own. Assets: no retained span contains a figure environment, a graphics-inclusion command, a PDF-page inclusion, a table or a TikZ picture; no image file is copied into this package, the package declares no asset dependency and no image file is redistributed with it. Licence: version arXiv:2312.07502v3 of this article is distributed under the Creative Commons Attribution 4.0 International licence, as recorded in the arXiv licence metadata for this exact version and shown on the abstract page https://arxiv.org/abs/2312.07502v3. A full rights notice is delivered at licenses/arxiv-2312.07502-CC-BY-4.0.txt. Characters: every retained excerpt is pure ASCII, so the delivered engine, XeLaTeX with its default Latin Modern text font, reports no missing character.
\begin{equation}  \label{incomplete gamma inequality}
\quad\left(1-e^{-s_a x}\right)^a<\frac{\gamma(a, x)}{\Gamma(a)}<\left(1-e^{-r_a x}\right)^a, \quad 0 \leq x<\infty, \quad a>0, \quad a \neq 1,
\end{equation}

\begin{equation}
\lim_{x\to\infty}\frac{\gamma(x+1, x)}{\Gamma(x+1)} = 1/2. \label{incomplete gamma go to 1}
\end{equation}

\begin{lemma}
Let $\{a_n\}, \{b_n\}>0$ be sequences such that  $a_n=O(1)$. Then, we have, as $n\to\infty$,
\begin{displaymath}
    \int_{0}^{a_n} x^{b_n-1} \exp(-x) dx \asymp a_n^{b_n}/b_n.
\end{displaymath}
\label{incomplete integrate}
\end{lemma}

\begin{lemma} \label{CH spectral density tail behavior}
Suppose $\beta^2=O(1)$ and $\alpha>d/2+1$, for $\alpha$ fixed or tending to infinity as $n \to \infty$. Then,
$$
(1+\lambda^2)^{v+d/2} m^{\alpha,\beta}_{CH}(\lambda) \gtrsim \left\{\begin{array}{lll}
\frac{\Gamma(\alpha-d/2)\beta^{d}}{\Gamma(\alpha)}(1+\lambda^{2})^{v+d/2},  & \text { if } \lambda^2 \beta^{2} \leq 1, \\
\frac{\beta^{-2v}\Gamma(\alpha-d/2)}{\Gamma(\alpha)}\frac{\alpha}{e^\alpha}, & \text { if } 1 < \lambda^2 \beta^{2} < \alpha+v-1,\\
\frac{\beta^{-2v}\Gamma(\alpha+v)}{\Gamma(\alpha)},& \text { if } \lambda^2 \beta^{2} \geq  \alpha+v-1.
\end{array}\right.
$$
\end{lemma}

\begin{proof}[Proof of Lemma \ref{CH spectral density tail behavior}] 
We only prove the $\alpha \to \infty$ case. The fixed $\alpha$ case can be proved by the same method. Let $ h=\beta^2/(2\phi^2)$. Then,
\begin{equation}\label{spectral density expansion}
\begin{split}
&(1+\lambda^2)^{(v+d/2)}m^{\alpha,\beta}_{CH}(\lambda) \\
=&  \frac{\sigma^2 2^{v-\alpha} v^v \beta^{2\alpha}}{\pi^{d/2}\Gamma(\alpha)}\int_0^{\infty} \left(\frac{1+\lambda^2}{2v\phi^{-2}+\lambda^2}\right)^{v+d/2}  \phi^{-2(v+\alpha+1)} \exp{(-\beta^2/(2\phi^2))d\phi^2}  \\
=&  \frac{\sigma^2  {(4v)}^v \beta^{-2v}}{\pi^{d/2}\Gamma(\alpha)}\int_0^{\infty} \left(\frac{1+\lambda^{2}}{4vh \beta^{-2}+\lambda^{2}}\right)^{v+d/2}  h^{v+\alpha-1} \exp(-h) dh\\
=&  \frac{\sigma^2  {(4v)}^v \beta^{-2v}}{\pi^{d/2}\Gamma(\alpha)} \left[\int_0^{\lambda^2\beta^2} \left(\frac{1+\lambda^{2}}{4vh \beta^{-2}+\lambda^{2}}\right)^{v+d/2}  h^{v+\alpha-1} \exp(-h) dh\right.
\\
&+\left.\int_{\lambda^2\beta^2}^{\infty} \left(\frac{1+\lambda^{2}}{4vh \beta^{-2}
+\lambda^{2}}\right)^{v+d/2}  h^{v+\alpha-1} \exp(-h) dh \right]\\
\asymp& \frac{\beta^{-2v}}{\Gamma(\alpha)} \left[\int_0^{\lambda^2\beta^2} \left(\frac{1+\lambda^{2}}{4vh \beta^{-2}+\lambda^{2}}\right)^{v+d/2}  h^{v+\alpha-1} \exp(-h) dh\right.
\\
&+\left.[(1+\lambda^{2})\beta^{2}]^{v+d/2}\int_{\lambda^2\beta^2}^{\infty} h^{\alpha-d/2-1} \exp(-h) dh \right].
\end{split}
\end{equation}
When $\lambda^2 \beta^{2}\leq 1$, by Lemma \ref{incomplete integrate}:
\begin{displaymath}
\begin{split}
&(1+\lambda^2)^{(v+d/2)}m^{\alpha,\beta}_{CH}(\lambda) \\
&\gtrsim \frac{\beta^{-2v}}{\Gamma(\alpha)} [(1+\lambda^{2})\beta^{2}]^{v+d/2}\int_{\lambda^2\beta^2}^{\infty} h^{\alpha-d/2-1} \exp(-h) dh \\
&\asymp \frac{\beta^{-2v}}{\Gamma(\alpha)}[(1+\lambda^{2})\beta^{2}]^{v+d/2} \Gamma(\alpha-d/2)\\
&\asymp \frac{\Gamma(\alpha-d/2)\beta^{d}}{\Gamma(\alpha)}(1+\lambda^{2})^{v+d/2}. 
\end{split}
\end{displaymath}
When $ \lambda^2 \beta^{2} \geq \alpha+v-1 $, we have, by (\ref{incomplete gamma go to 1}) and (\ref{spectral density expansion}):
%\begin{footnotesize}
\begin{displaymath}
\begin{split}
&(1+\lambda^2)^{(v+d/2)}m^{\alpha,\beta}_{CH}(\lambda) \\
&\asymp \frac{\beta^{-2v}}{\Gamma(\alpha)} \left[\int_0^{\lambda^2\beta^2}  h^{v+\alpha-1} \exp(-h) dh
+[(1+\lambda^{2})\beta^{2}]^{v+d/2}\int_{\lambda^2\beta^2}^{\infty} h^{\alpha-d/2-1} \exp(-h) dh \right]\\
&\gtrsim \frac{\beta^{-2v}\Gamma(\alpha+v)}{\Gamma(\alpha)}.
\end{split}
\end{displaymath}
When $ 1\leq \lambda^2 \beta^{2} \leq \alpha+v-1 $, we have, by (\ref{incomplete gamma inequality}), (\ref{spectral density expansion})  and $e^x \ge 1+x$:
\begin{displaymath}
\begin{split}
&(1+\lambda^2)^{(v+d/2)}m^{\alpha,\beta}_{CH}(\lambda) \\
\asymp& \frac{\beta^{-2v}}{\Gamma(\alpha)} \left[\int_0^{\lambda^2\beta^2}  h^{v+\alpha-1} \exp(-h) dh
+[(1+\lambda^{2})\beta^{2}]^{v+d/2}\int_{\lambda^2\beta^2}^{\infty} h^{\alpha-d/2-1} \exp(-h) dh \right]\\
\gtrsim& \frac{\beta^{-2v}}{\Gamma(\alpha)}[(1+\lambda^{2})\beta^{2}]^{v+d/2}\int_{\lambda^2\beta^2}^{\infty} h^{\alpha-d/2-1} \exp(-h) dh\\
\gtrsim& \frac{\beta^{-2v}}{\Gamma(\alpha)}\int_{ \alpha+v-1}^{\infty} h^{\alpha-d/2-1} \exp(-h) dh\\
\geq& \frac{\beta^{-2v}\Gamma(\alpha-d/2)}{\Gamma(\alpha)}[1-(1-e^{- \alpha-v+1})^{\alpha-d/2}]\\
%=& \frac{\beta^{-2v}\Gamma(\alpha-d/2)}{\Gamma(\alpha)}[1-(1-e^{- \alpha-v+1})^{e^{ \alpha+v-1} \frac{\alpha-d/2}{e^{ \alpha+v-1}}}]\\
\gtrsim&\frac{\beta^{-2v}\Gamma(\alpha-d/2)}{\Gamma(\alpha)}\left[1-\exp\left(-{ \frac{\alpha-d/2}{e^{ \alpha+v-1}}}\right)\right]\\
\gtrsim& \frac{\beta^{-2v}\Gamma(\alpha-d/2)}{\Gamma(\alpha)}{ \frac{\alpha-d/2}{e^{ \alpha+v-1}}}\\
\gtrsim&\frac{\beta^{-2v}\Gamma(\alpha-d/2)}{\Gamma(\alpha)}\frac{\alpha}{e^\alpha}. 
\end{split}
\end{displaymath}
%\end{footnotesize}
\end{proof}

\end{document}
